\section{Back-office : un effet \enquote{boîte noire}}
Si la bibliothèque numérique fait preuve de robustesse, sa limite majeure réside en coulisses: l'architecture du back-office. Face à un système fermé et dépourvu de documentation, les équipes métiers se retrouvent privées d'autonomie et de visibilité sur l'ensemble de la chaîne de traitement, ce qui entrave considérablement les cadences de production. 

Chaque tâche interne devient plus complexe notamment lorsque l'un d'eux repère une erreur syntaxique. Par exemple, ils peuvent la corriger mais ensuite ils doivent relancer toute la procédure de publication, ce qui prend du temps là, où une petite correction peut prendre entre une et deux minutes à être corrigée, et devrait être visible presque intantanément sur le portail de la bibliothèque \footannexe{\nameref{annexe_A_rapport_audit}.}. Ce n'est actuellement pas le cas pour la Bibliothèque Cujas qui se trouve contrainte de relancer son protocole de publication afin d'afficher l'erreur rectifiée. En effet, une fois l'erreur rectifiée, la personne en charge de la publication doit à nouveau rouvrir le \ac{PDF}, scellé par un mot de passe, modifier sa date de création, le verrouiller à nouveau, avant de le mettre en ligne\footannexe{\nameref{annexe_A_rapport_audit}.}. À chaque étape, ils doivent attendre que les traitements de nuit valident les étapes précédentes avant de pouvoir poursuivre les modifications, ajoutant une lourdeur à la procédure de travail \footannexe{\nameref{annexe_A_rapport_audit}.}. 

Le second problème vient de \ac{XTF} lui-même qui n'est plus maintenu. Habituellement, les inconvénients quotidiens ou difficultés rencontrées par les experts métiers étaient tous palliés par le responsable du \ac{DIT} \footannexe{\nameref{annexe_A_rapport_audit}.}. Toutefois, c'est loin d'être une bonne chose, car \enquote{tous les services de la bibliothèque numérique actuelle qu'ils soient interne ou externe reposent uniquement sur les compétences d'une seule et même personne}\footannexe{\nameref{annexe_A_rapport_audit}.}. Aujourd'hui, la bibliothèque numérique actuelle survit uniquement grâce aux compétences informatiques du responsable qui a su trouver des solutions comme l'ajout d'une visionneuse \ac{JPEG} permettant de consulter les ouvrages volumineux. 

Etant loin d'être un cas isolé, la situation de la Bibliothèque Cujas illustre bien la pensée de l'époque des institutions patrimoniales, ainsi que le concept du deuxième âge théorisé par Emmanuelle Bermès et Gautier Poupeau. Ce deuxième âge s'étend des années 1990 à 2010. Deux événements majeurs marquent cette période: l'irruption du Web au milieu des années 1990 et la numérisation de masse insufflée par Google à partir des années 2000 \footcite{bermesTroisAgesNumerique}. Cette époque correspond aux développements des projets de numérisation de la Bibliothèque Cujas qui s'est structurée selon les pratiques de travail de l'époque. Par ailleurs, c'est à ce moment-là qu'émerge les premières bibliothèques numériques \footcite{bermesTroisAgesNumerique}. 

Entre 1990 et 2010, les institutions patrimoniales ont une vision différente de la conception d'un document numérique. À l'épopque, les bibliothèques voient le document numérique comme une conversion du document physique. Jean-Michel Salaün théorise d'ailleurs cette idée qu'il définit \enquote{comme un ensemble de données organisées selon une structure fiable associées à des règles de mise en forme permettant une lisibilité partagée entre son concepteur et ses lecteurs} \footcite{pedauqueDocumentLumiereNumerique2006}. Avec l'émergence du web, de nouveaux standards se développent, utiles pour structurer les notices bibliographiques. C'est l'arrivée massif du format \ac{XML} qui permet le triomphe \enquote{de la pensée hiérarchique et documentaire}, modélisant les documents sous forme d'arborescences \footcite{bermesTroisAgesNumerique}. Le format \ac{XML} est alors pensé pour permettre la diffusion des données avec les autres bibliothèques sauf que l'arrivée des ordinateurs personnels freine paradoxalement l'innovation. Christine Ducourtieux l'exprime d'ailleurs ainsi : \enquote{l'arrivée des systèmes d'exploitation Windows et MacIntosh et de la bureautique ont paradoxalement pour effet de freiner l'élan d'innovation technique dans les méthodes, remplacé progressivement par des patriques discrètes : une appropriation de l'outil numérique plus individuelle, qui n'est plus subordonnée à l'histoire quantitative} \footcite{bermesTroisAgesNumerique}. Le développement de la technologie et l'arrivée de l'Internet en 1989, entraîne à nouveau une forme de cloisonnement avec un retour à l'individualisation des pratiques des chercheurs. 

En effet, l'arrivée des ordinateurs personnels bouleversent les pratiques créant une \enquote{séparation nette} entre l'édition scientifique et la numérisation de masse. Emmanuelle Bermès justifie cette hypothèse en expliquant que \enquote{l'édition scientifique numérique, qui reste le domaine de la recherche et peut s'adapter aux spécificités des corpus, et la numérisation, qui vise des masses documentaires plus importantes et traitées de manière homogène} \footcite{bermesTroisAgesNumerique}. Loin de favoriser la collaboration, cette évolution informatique tend à réisoler les pratiques de travail. Les bibliothèques numériques risquent alors d'être réduites à de simples entrepôts de fichiers \ac{PDF}, déconnectées des besoins de valorisation et du travail des chercheurs.

On observe aussi une parcellisation industrielle des tâches en bibliothèque. À cette époque, les tâches sont énormément segmentisées. Chaque personne a uniquement accès à ce qu'il doit faire, considérant qu'il n'est pas utile pour le bon déroulement de son travail d'avoir accès à l'intégralité du processus de production. La Bibliothèque Cujas est héritière de cette manière de travail mais c'est loin d'être un cas isolé car d'autres institutions patrimoniales comme la \ac{BNF} ont également rencontré ce problème. Michel Richard écrit sur ce sujet :\enquote{l'équipe responsable de la réalisation de ce programme comprend quinze personnes à plein temps rattachées sur le site d'Ivry. elle a en charge l'ensemble des tâches nécessaires à la constitution de la collection (travail bibliographique, contacts avec les universités et l'édition, relations avec la recherche, gestion du fonds électronique, commandes aux libraires, négocations avec les ayants droit, etc.) et au suivi de production (choix des prestataires, tests, contrôle qualité, gestion des stocks, normalisation, etc.)} \footcite{bermesTroisAgesNumerique}. Cela alourdit énormément le travail des bibliothécaires car c'est uniquement le \enquote{résultat final} qui est mutualisé et non le reste.  

La Bibliothèque Cujas fait d'ailleurs face à cette contrainte en disant clairement que le système actuel survit uniquement grâce aux compétences informatiques du responsable du \ac{DIT}. En dehors de lui, personne n'est à même à Cujas de pouvoir régler le problème entraînant ainsi une centralisation des compétences en une seule personne mettant en péril la pérennité de la bibliothèque numérique dans le futur. Les données deviennent alors vulnérables dans un monde où avec l'arrivée du Web sémantique, on réfléchit à ouvrir d'avantage les données afin de permettre une libre circulation de celle-ci.